ffgraph is a C++ class library to create, manipulate and display
directed graphs. It was developed in the context of the
\meta{}\cite{meta:95} project to visualize the various data structures
used in this project.

There already exist many graph-tools, so the question arises why to
write a new one instead of using an existing one. Many of the existing
ones are programs like GraphEd\cite{GraphEd:Himsolt-90b} or the VCG
tool\cite{sander:94} so you can use them to display your data, but you
have to implement a graph data structure in your program anyway. If you
use external programs to display your graphs it is often impossible to
dynamically change the graph in your program and update the display, or
to manipulate it graphically and pass the update back to your
program. On the other hand existing graph libraries, like
LEDA\cite{naeher:94} offer a powerful graph data structure, but have
less powerful means for displaying the data. Other tools are specialised
on a small subset of graphs, like dataflow graphs etc., and are thus not
usable for us. So we decided to write a graph library on our own, which
offers you a graph data structure and powerful means to display it.

Graphs consist of a set of nodes and a set of
edges, where edges have exactly one target and one source. Access
to the library is provided through the classes described in this
paper. Graphs can be loaded and saved in a slightly extended
sgraph\cite{Sgraph:Himsolt-93a} file format\footnote{Compare
  section \vref{sgraph_file_format}}. The main purpose of this
library is not speed or the ability to handle huge graphs, but to
give an easy to use and save interface, the possibility to bind
user specific data to graphs, nodes, and edges and to provide a
powerful and flexible mechanism to display graphs. 

The library allows the user to bind semantic labels to graphs, nodes,
and edges. Labels are C++ classes which may contain any user specific
methods or data and get informed of any changes to the structure of
the graph. Some labels are included in the distribution to give an
example on how to use labels\footnote{E.g. there is a label
implementing nested graphs included in the distribution.}.  

Special graph labels which are subclasses of {\bf FFGraphWindow} may be
assigned to a graph in order to display it in up to 3 dimensions. The
user may provide his/her own window labels to display the graph the way
it best suits his/her demands. The distribution contains window labels
which display 2 or 3 dimensional graphs using the tcl/tk\cite{Tcl:Oust}
widget library for X11 systems\footnote{See classes {\bf
FFTclBasicWindow}, {\bf FFTclWindow} and {\bf FFTclSubWindow}}.

To display the contents of labels, inspectors are used. The class {\bf
FFTclWindow} contains an inspector that allows one to inspect the
content of a label or to invoke the inspector of a label if the label
has an inspector on its own.

A special interface is provided for labels that implement layout
algorithms. Two layout labels implementing a 3 dimensional
spring-embedder and a sugiyama layout algorithm\footnote{Written by Falk
Schreiber (schreibe{@}fmi.uni-passau.de)} are included in the
distribution.

A graph editor is included in the distribution to give an example
of how to use the library.

There are several standard template container classes like sets,
maps, etc. included in the distribution. These classes are naively
modified and debugged versions of the libg++ container classes
and are to be replaced soon by more efficient implementations and in some
cases by better types of containers. For more information on how
to use the container classes see the libg++ documentation.

The library also contains a slightly modified version of a Xpm
extension to tk 4.0 by Mark Crimmins (markcrim{@}umich.edu) which
is based on a work by Harald Albrecht
(albrecht{@}igpm.rwth-aachen.de).

\newpage
